% Part: many-valued-logic
% Chapter: three-valued-logics
% Section: introduction

\documentclass[../../../include/open-logic-section]{subfiles}

\begin{document}

\olfileid{mvl}{thr}{int}

\olsection{پېژندنه}

که موږ له~$\True$ او~$\False$ سره يوازې يو بل قيمت~$\Undef$
ورزيات کړو، د درې ارزښته منطق د ماتريس د قيمتونو سټ ترلاسه کوو۔ که څه هم د رښتياوالي
يوازې يو قيمت ورزيات شوے، د~$\lnot$، $\land$، $\lor$ او
$\lif$ د صدق تابعو د تعريفولو امکانونه ډېر دي۔ بيا يو منطق د دغو
صدق تابعو هره ګډوله کارولے شي۔ په دې باب کښې به د ټاکلو قيمتونو دوه انتخابونه څېړو: يا يوازې
$\True$ ټاکل شوے وي، يا~$\True$ او~$\Undef$ دواړه ټاکل شوي وي۔

دلته موږ ځينې تر ټولو پېژندل شوي درې ارزښته منطقونه، د هغو
انګېزې او ځينې خاصيتونه وړاندې کوو۔

\end{document}
